% ============================================================
%  EXPERIMENTS SECTION — MO-QAS PAPER (code-aligned + PDF figures)
% ------------------------------------------------------------
%  Preamble (main document):
%    \usepackage{amsmath,booktabs,graphicx,hyperref}
%  If you \input{} this file from a subfolder, either:
%    \graphicspath{{figures/}}   % when main.tex is in repo root, or
%    \graphicspath{{../figures/}}
%  Alternatively redefine \FigPath once before \input:
%    \renewcommand{\FigPath}{../figures/}
% ------------------------------------------------------------
%  PDFs (generated by scripts under figures/):
%    gradient_scaling.pdf, noise_sweep.pdf, fidelity_heatmap.pdf
% ============================================================

\providecommand{\FigPath}{figures/}
\providecommand{\path}[1]{\texttt{#1}}
\newcommand{\Incfig}[2][width=\linewidth]{%
  \includegraphics[#1]{\FigPath#2}}

\section{Experiments}\label{sec:experiments}

This section evaluates three \emph{complementary but distinct} threads.
We state the separation so readers do not conflate them.

\begin{enumerate}
  \item[\textbf{(A)}] \textbf{Surrogate-driven MO-QAS.}
        A Graph Transformer is trained on simulator-labeled architectures, then
        NSGA-II uses \emph{only} surrogate forward passes as fitness---no noisy
        density-matrix simulation in the inner loop. Outputs:
        \path{experiments/moqas_pareto_results.json}.

  \item[\textbf{(B)}] \textbf{Noise-channel fidelity (fixed templates).}
        Two \emph{hand-specified} PennyLane circuits (depth-6 vs.\ depth-2,
        \textbf{linear} CNOT ladders) on \path{default.mixed}. This is
        \emph{not} an exhaustive re-evaluation of every Pareto genome from~(A).

  \item[\textbf{(C)}] \textbf{Gradient-variance scaling (fixed synthetic families).}
        Three fixed families at increasing~$n$; \emph{not} chromosomes from~(A).
        Supports the trainability narrative; it does not prove scaling for
        evolved architectures.
\end{enumerate}

\noindent
Raw outputs are JSON under the repo root and \path{experiments/}.
Table~\ref{tab:figure_map} maps figures to scripts; Section~\ref{ssec:repro}
lists commands.

% ------------------------------------------------------------------
\subsection{Overview of Reported Figures}\label{ssec:figure_overview}
% ------------------------------------------------------------------

\paragraph{Figure~\ref{fig:gradient_scaling} (gradient variance vs.\ $n$).}
Horizontal axis: qubit count $n\in\{4,6,8,10,12\}$.
Vertical axis: mean gradient variance (log scale in the plot).
Three \textbf{fixed} families from \path{experiments/gradient_scaling_standalone.py}:
\textsf{Global} --- depth~10, \textbf{all-to-all} CNOTs, \textbf{global} cost
$\langle Z^{\otimes n}\rangle$;
\textsf{Local} --- depth~10, \textbf{all-to-all} CNOTs (nested loops over all
pairs), \textbf{local} cost $\langle Z_0\rangle$;
\textsf{NAS} (MO-QAS-style) --- depth~6, \textbf{linear} CNOT chain, local
$\langle Z_0\rangle$.
Each point uses up to \texttt{n\_samples=40} Gaussian parameter draws
(scale~$0.1$); the script averages coordinate-wise gradient variances across
draws (see \texttt{compute\_gradient\_variance} in that file).
Data: \path{experiments/gradient_scaling_results.json}; script:
\path{figures/figure1_gradient_scaling.py}.
\textbf{Takeaway:} empirical scaling of gradient variance with~$n$ for these
templates; fitted slopes (legend) come from log--log regression on the saved
series (Table~\ref{tab:grad_exponents}). The \textsf{Global} curve has the
\emph{largest} variance in this experiment at small~$n$ and the least negative
fitted exponent (slowest decay with~$n$ among the three); \textsf{Local} decays
fastest (most negative~$\alpha$).

\paragraph{Figure~\ref{fig:noise_sweep} (fidelity vs.\ noise probability).}
Three panels (depolarizing, amplitude damping, phase damping).
$x$-axis: noise probability in percent; $y$-axis: fidelity
$F\in[0,1]$ from \texttt{qml.math.fidelity}$(\rho_{\mathrm{ideal}},\rho_{\mathrm{noisy}})$
with both arguments as density matrices on \path{default.mixed}.
Curves: depth-6 baseline vs.\ depth-2 knee\emph{-style} template, $n=4$ block
in \path{multi_noise_results_seed42.json}; script:
\path{figures/figure2_noise_sweep.py}.
\textbf{Takeaway:} the shallower template retains higher fidelity as $p$ grows for
all three channels in this run; depolarizing panel annotates the $F_2/F_6$ ratio
at $10\%$ noise ($2.55$, matching Table~\ref{tab:fidelity_10pct}).

\paragraph{Figure~\ref{fig:fidelity_heatmap}.}
Heatmap of the \textbf{same} JSON: fidelity vs.\ noise level and depth template,
one block per channel; script \path{figures/figure3_fidelity_heatmap.py}.
\textbf{Takeaway:} compact view of the sweep behind Figure~\ref{fig:noise_sweep}
and Table~\ref{tab:fidelity_10pct}.

% ------------------------------------------------------------------
\subsection{Experiment~(A): MO-QAS Protocol}\label{ssec:moqas}
% ------------------------------------------------------------------

\subsubsection{Dataset and Surrogate}

Architectures are encoded as a \textbf{bipartite graph} via
\path{ImprovedBipartiteGraphBuilder} (architecture and constraint nodes),
\emph{not} a per\hyp gate DAG.
The training pickle \path{paper_dataset.pkl} stores tuples produced by
evaluating random sampled architectures on the \textbf{digits} classification
task with the project simulator; labels supervise the surrogate heads.

\path{GraphTransformerPredictor} applies graph-attention layers; readouts feed
three heads: accuracy, depth (raw network output $\hat d_{\mathrm{raw}}$ mapped to
$\texttt{depth\_est}=\max(0,\hat d_{\mathrm{raw}})\times 40$ in NSGA\hyp II
scoring), and trainability in $[0,1)$ (clamped before $\hat V$).
Let $t$ be trainability; the code uses
$\hat V=\max(10^{-8},\,1-t)$ and the third objective term $-\log \hat V$.
The factor \textbf{40} is a fixed scaling used to keep depth estimates in a
comparable numeric range during optimization diagnostics; dominance direction is
unchanged, and knee selection later normalizes objectives within the front.

\subsubsection{Objectives and Dominance}

NSGA-II maximizes three scalar scores $(f_1,f_2,f_3)$ on every candidate
(Pareto dominance uses component-wise $\geq$ with at least one strict
inequality):
\begin{equation}\label{eq:moqas_objectives}
  (f_1,\,f_2,\,f_3)
  =
  \bigl(
    \hat a,\;
    -\texttt{depth\_est},\;
    -\log \hat V
  \bigr).
\end{equation}
Higher $f_1$ is better accuracy; higher $f_2$ means smaller \texttt{depth\_est};
higher $f_3$ means larger trainability (smaller~$\hat V$). No density-matrix
noise simulation runs inside the evolutionary loop.

\subsubsection{Hyperparameters}

Values are fixed in \path{run_moqas_pipeline.py}.

\begin{table}[t]
  \centering
  \caption{MO-QAS pipeline hyperparameters (\path{run_moqas_pipeline.py}).}
  \label{tab:moqas_hyper}
  \begin{tabular}{lcc}
    \toprule
    Parameter & \texttt{--quick} & Full \\
    \midrule
    Dataset circuits & 48 & 120 \\
    Batch size & 8 & 8 \\
    Surrogate epochs & 25 & 60 \\
    NSGA-II population & 12 & 20 \\
    NSGA-II generations & 8 & 15 \\
    \bottomrule
  \end{tabular}
\end{table}

\subsubsection{Knee-Point Selection (matches \path{quantum_nas.py})}

Let $\mathcal{F}$ index the final non-dominated front.
For $i\in\mathcal{F}$, let $a_i$ be surrogate accuracy, $d_i=\texttt{depth\_est}$,
and $g_i=-\log \hat V_i$.
Min-max normalize \emph{within the front} (code uses a range or $1.0$ if
degenerate):
\[
  \tilde a_i = \frac{a_i - \min_{j\in\mathcal F}a_j}{\max_{j\in\mathcal F}a_j - \min_{j\in\mathcal F}a_j + \epsilon_a},\quad
  \tilde d_i = \frac{\max_{j\in\mathcal F}d_j - d_i}{\max_{j\in\mathcal F}d_j - \min_{j\in\mathcal F}d_j + \epsilon_d},
\]
\[
  \tilde g_i = \frac{g_i - \min_{j\in\mathcal F}g_j}{\max_{j\in\mathcal F}g_j - \min_{j\in\mathcal F}g_j + \epsilon_g}.
\]
Here $\epsilon_a,\epsilon_d,\epsilon_g$ denote numerical guards for zero-range
cases; implementation-equivalently, each denominator uses the observed range when
positive and $1.0$ when degenerate.
The knee is $\arg\max_{i\in\mathcal F}\min(\tilde a_i,\tilde d_i,\tilde g_i)$
(\texttt{\_find\_knee\_moqas}).

\subsubsection{Outputs (illustrative: \path{moqas_pareto_results.json} from a full run)}

\begin{itemize}
  \item Non-dominated solutions on surrogate objectives: \textbf{20} points
        (one saved run).
  \item Knee (surrogate): $\hat a\approx 0.919$;
        $\texttt{depth\_est}\approx 8.74$ (this is the scaled depth used in
        $(f_1,f_2,f_3)$; it is \emph{not} the logical layer count);
        $-\log\hat V\approx 0.952$;
        decoded architecture: logical depth~\textbf{2}, \textbf{10}~qubits, RY,
        hardware-efficient parametrization (the JSON may still list e.g.\
        \texttt{circular} entanglement in the chromosome).
\end{itemize}
\noindent
The currently checked-in MO-QAS JSON used for bridge tables is produced by the
\textbf{full} configuration (\texttt{dataset\_samples=120},
\texttt{predictor\_epochs=60}, \texttt{pop\_size=20},
\texttt{generations=15}); bridge tables below are regenerated from this file.

\paragraph{Caveat.}
Surrogate labels come from the dataset simulator setup, \textbf{not} from the
mixed-state channels in Experiment~(B).

% ------------------------------------------------------------------
\subsection{Experiment~(B): Noise Robustness (Fixed Templates)}\label{ssec:noise}
% ------------------------------------------------------------------

\subsubsection{Protocol (\path{multi_noise_validation_verified.py})}

\textbf{Circuits are explicit gate sequences} (e.g.\ \texttt{RY}, \texttt{RX},
\texttt{RZ}, \texttt{CNOT}), not a template class such as
\texttt{StronglyEntanglingLayers}:
\begin{itemize}
  \item \textbf{Depth~6:} one round of $RY(\pi x_i)$ per qubit with
        $x_i\equiv 0.5$; then six layers, each with noise on one-qubit wires,
        $RX(\pi/4)$ per qubit, linear CNOTs with noise on CNOT wires.
  \item \textbf{Depth~2:} two layers; each layer repeats the $RY$ encoding,
        noise on one-qubit wires, $RX(\pi/4)$ and $RZ(\pi/4)$ per qubit,
        linear CNOTs with noise on CNOT wires (data re-uploading).
\end{itemize}
Parameters are \textbf{fixed} ($\pi/4$, inputs $0.5$), not trained.
Noise strengths $p\in\{0,0.02,0.05,0.10\}$; channels:
depolarizing, amplitude damping, phase damping.
Primary tables and figures use $n=4$; the JSON also stores $n=5$.

\begin{table}[t]
  \centering
  \caption{Noise-fidelity experiment configuration.}
  \label{tab:noise_config}
  \begin{tabular}{ll}
    \toprule
    Item & Value \\
    \midrule
    Backend & \path{default.mixed} \\
    Qubits (main table/figs.) & 4 (5 in JSON also) \\
    Depths & 6 vs.\ 2 (linear CNOTs) \\
    Channels & Depolarizing, amplitude damping, phase damping \\
    $p$ & $\{0,0.02,0.05,0.10\}$ \\
    Primary seed & 42 \\
    Multi-seed & \texttt{verify} $\rightarrow$ $\{42,123,456\}$ \\
    \bottomrule
  \end{tabular}
\end{table}

\subsubsection{Results at $p=0.10$ ($n=4$, seed~42)}

\begin{table}[t]
  \centering
  \caption{State fidelity at $p=10\%$ (from \path{multi_noise_results_seed42.json}, \texttt{"4"} block).}
  \label{tab:fidelity_10pct}
  \begin{tabular}{lccc}
    \toprule
    Noise channel & Depth-6 $F$ & Depth-2 $F$ & Ratio (D2/D6) \\
    \midrule
    Depolarizing      & 0.0833 & 0.2121 & 2.55 \\
    Amplitude damping & 0.2514 & 0.3986 & 1.59 \\
    Phase damping     & 0.2514 & 0.6223 & 2.48 \\
    \bottomrule
  \end{tabular}
\end{table}

\begin{figure*}[t]
  \centering
  \Incfig[width=\textwidth]{noise_sweep.pdf}
  \caption{State fidelity vs.\ noise probability (\%) for depth-6 baseline
    (red) and depth-2 knee-style template (green), $n=4$, seed~42.
    Three channels: depolarizing, amplitude damping ($T_1$), phase damping ($T_2$).
    Source: \path{figures/noise_sweep.pdf} from \path{figure2_noise_sweep.py}.}
  \label{fig:noise_sweep}
\end{figure*}

\begin{figure*}[t]
  \centering
  \Incfig[width=\textwidth]{fidelity_heatmap.pdf}
  \caption{Fidelity heatmap for the same data as Figure~\ref{fig:noise_sweep}:
    rows depth-6 baseline vs.\ MO-QAS knee-style template; columns
    $p\in\{0,2,5,10\}\%$; one panel per noise channel.
    Source: \path{figures/fidelity_heatmap.pdf} from \path{figure3_fidelity_heatmap.py}.}
  \label{fig:fidelity_heatmap}
\end{figure*}

\noindent
\texttt{python multi\_noise\_validation\_verified.py verify} writes
\path{multi_noise_results_3seeds.json}.

\subsubsection{Direct A $\rightarrow$ B Bridge Check (three Pareto representatives)}

We decode representative Pareto architectures from
\path{experiments/moqas_pareto_results.json}
(\textbf{full} config in that file), then simulate each under the same
mixed-state channels at $p=0.10$ on \path{default.mixed}. Script:
\path{experiments/bridge_knee_noise_eval.py}. Output JSON:
\path{experiments/bridge_pareto_noise_results.json}.
Representative qubit counts differ because qubit count is an architecture
decision variable in the MO-QAS search space; therefore, cross-row ratio
comparisons should be interpreted with care. For tractable mixed-state runtime,
the bridge script caps representatives at $n\le 10$.

\begin{table*}[t]
  \centering
  \caption{Direct Pareto re-evaluation from Experiment~(A) at $p=10\%$
           (\path{experiments/bridge_pareto_noise_results.json}).}
  \label{tab:bridge_pareto_10pct}
  \footnotesize
  \setlength{\tabcolsep}{4pt}
  \begin{tabular}{lcccccccc}
    \toprule
    Representative & $n$ &
    Depol. $F_{\text{arch}}$ & Depol. $F_{D6}$ &
    Amp. $F_{\text{arch}}$ & Amp. $F_{D6}$ &
    Phase $F_{\text{arch}}$ & Phase $F_{D6}$ \\
    \midrule
    Knee (circular) & 10 & 0.0094 & 0.0013 & 0.0860 & 0.0180 & 0.2444 & 0.0180 \\
    High-accuracy non-knee unique (circular) & 8 & 0.0063 & 0.0052 & 0.0471 & 0.0439 & 0.1824 & 0.0439 \\
    Low-depth remaining unique (linear) & 10 & 0.0139 & 0.0013 & 0.0984 & 0.0180 & 0.2765 & 0.0180 \\
    \bottomrule
  \end{tabular}
\end{table*}

\noindent
At $p=10\%$, architecture-to-baseline ratios computed from these values are:
for knee ($n=10$): depol.\ $7.11$, amp.\ $4.77$, phase $13.54$; for the other
two representatives ($n=8$ and $n=10$): depol.\ $1.21/10.58$, amp.\ $1.07/5.45$,
phase $4.16/15.32$.

\subsubsection{What~(B) Does \emph{Not} Show}

\begin{itemize}
  \item Three decoded Pareto representatives are simulated, but this is still not
        an exhaustive re-evaluation of all Pareto members.
  \item No hardware-specific noise.
  \item Fidelity is structural (fixed $\theta$), not best-trained accuracy under noise.
\end{itemize}

% ------------------------------------------------------------------
\subsection{Experiment~(C): Gradient-Variance Scaling}\label{ssec:grad}
% ------------------------------------------------------------------

Protocol matches \path{experiments/gradient_scaling_standalone.py}:
for each $(n,\text{family})$, sample parameters, parameter-shift gradient of the
full vector, retain valid samples, then mean of coordinate-wise variances
across samples (\texttt{np.mean(np.var(...,\,axis=0))}).

\paragraph{Observables.}
Global cost: tensor product $Z^{\otimes n}$ (expectation via chained
\texttt{PauliZ}); local cost: $\langle Z_0\rangle$.

\begin{table}[t]
  \centering
  \caption{Log-log regression exponents $\alpha$ in
           $\mathrm{Var}\propto n^{\alpha}$ from
           \path{experiments/gradient_scaling_results.json}
           (values rounded to two decimals; re-run script to refresh).}
  \label{tab:grad_exponents}
  \begin{tabular}{lc}
    \toprule
    Family & $\alpha$ \\
    \midrule
    Global (depth 10, full, global cost) & $-1.20$ \\
    Local (depth 10, full, local cost)  & $-1.97$ \\
    NAS (depth 6, linear, local cost)   & $-1.57$ \\
    \bottomrule
  \end{tabular}
\end{table}

\begin{figure*}[t]
  \centering
  \Incfig[width=\textwidth]{gradient_scaling.pdf}
  \caption{Mean gradient variance vs.\ qubit count for three fixed families
    (\textsf{Global}, \textsf{Local}, \textsf{NAS}/MO-QAS-style), shown at
    full width for legibility in two-column layout. Source:
    \path{figures/gradient_scaling.pdf} from \path{figures/figure1_gradient_scaling.py}.}
  \label{fig:gradient_scaling}
\end{figure*}

\noindent
More negative~$\alpha$ implies faster variance decay with~$n$ on this
log-log fit (finite grid $n\in\{4,6,8,10,12\}$, stochastic sampling).

\subsubsection{Direct A $\rightarrow$ C Bridge Check (knee architecture at fixed $n$)}

To partially connect evolved architectures to Experiment~(C), we compute the
same gradient-variance statistic for the decoded MO-QAS knee architecture at its
native qubit count (\path{experiments/bridge_knee_gradient_eval.py}).
Result (\path{experiments/bridge_knee_gradient_result.json}):
\[
\mathrm{Var}_{\mathrm{knee}}(n{=}10)=1.30\times 10^{-8},
\]
with reference NAS-template value at $n=10$ from
\path{experiments/gradient_scaling_results.json}:
\[
\mathrm{Var}_{\mathrm{NAS\ template}}(n{=}10)=8.40\times 10^{-4},
\]
ratio $\approx 1.55\times 10^{-5}$.
For one additional Pareto representative (low-depth remaining unique, linear,
same $n=10$), the measured variance is
\[
\mathrm{Var}_{\mathrm{low\mbox{-}depth}}(n{=}10)=3.53\times 10^{-5},
\]
which remains below the fixed NAS-template reference while being above the
knee-point value.
This is consistent with a structural hypothesis: the decoded knee circuit is
shallower (depth~2) than the NAS reference template used in the scaling sweep
(depth~6), and its learned structure (hardware-efficient with circular coupling)
differs from the fixed NAS template (linear chain), which changes effective
expressivity and gradient concentration. At $n=10$, both are already in a
low-variance regime; the bridge value indicates the decoded knee sits substantially
below the fixed template baseline. This observation is descriptive rather than
causal evidence; targeted ablations are needed to isolate which architectural
factor drives the reduction.

\paragraph{What~(C) Does \emph{Not} Show.}
No bootstrap CIs; no guarantee NSGA-II solutions match these templates;
$n\le 12$ only. The direct A$\rightarrow$C bridge above is a single-$n$ point,
not a scaling law over evolved architectures.

% ------------------------------------------------------------------
\subsection{Cross-Thread Validation Bridge (A $\rightarrow$ B/C)}\label{ssec:bridge}
% ------------------------------------------------------------------

This revision adds a bridge by directly re-evaluating three decoded Pareto
representatives under the Experiment~(B) noise protocol
(Table~\ref{tab:bridge_pareto_10pct}) and by adding a single-point A$\rightarrow$C
gradient statistic for the knee architecture.
However, Experiments~(B) and~(C) still mostly use fixed templates instead of a
representative Pareto subset. The next strengthening step is to evaluate multiple
Pareto points across multiple random seeds / full-config runs and to compute
gradient-variance statistics for those same architectures across $n$ where possible.

Interpret the current evidence as:
\textbf{(A)} surrogate-guided search behavior, \textbf{(B)} structural
noise-sensitivity of fixed circuits, and \textbf{(C)} scaling behavior of fixed
families, with an initial direct bridge rather than a full architecture-level
validation chain.
Full-config surrogate search is now integrated into the bridge tables; multi-seed
convergence diagnostics remain the main remaining upgrade for camera-ready
submission.

% ------------------------------------------------------------------
\subsection{Figure--Data--Script Map}\label{ssec:figure_map}
% ------------------------------------------------------------------

\begin{table}[t]
  \centering
  \caption{Figure provenance (run scripts from repo root). PDFs load from \texttt{figures/} unless you redefine \texttt{\textbackslash FigPath}.}
  \label{tab:figure_map}
  \begin{tabular}{lll}
    \toprule
    Figure & Script & Data \\
    \midrule
    \ref{fig:gradient_scaling}
      & \path{figures/figure1_gradient_scaling.py}
      & \path{experiments/gradient_scaling_results.json} \\
    \ref{fig:noise_sweep}
      & \path{figures/figure2_noise_sweep.py}
      & \path{multi_noise_results_seed42.json} \\
    \ref{fig:fidelity_heatmap}
      & \path{figures/figure3_fidelity_heatmap.py}
      & \path{multi_noise_results_seed42.json} \\
    \bottomrule
  \end{tabular}
\end{table}

\begin{table}[t]
  \centering
  \caption{Non-figure data provenance for bridge results.}
  \label{tab:nonfigure_provenance}
  \begin{tabular}{lll}
    \toprule
    Result & Script & Data \\
    \midrule
    A$\rightarrow$B Pareto bridge (Table~\ref{tab:bridge_pareto_10pct})
      & \path{experiments/bridge_knee_noise_eval.py}
      & \path{experiments/bridge_pareto_noise_results.json} \\
    A$\rightarrow$C knee gradient point
      & \path{experiments/bridge_knee_gradient_eval.py}
      & \path{experiments/bridge_knee_gradient_result.json} \\
    \bottomrule
  \end{tabular}
\end{table}

% ------------------------------------------------------------------
\subsection{Reproducibility}\label{ssec:repro}
% ------------------------------------------------------------------

Run from the repository root.

\begin{verbatim}
python run_moqas_pipeline.py --quick
python run_moqas_pipeline.py
# Optional: --skip-noise --skip-gradient --skip-figures

python multi_noise_validation_verified.py 42
python multi_noise_validation_verified.py verify
python experiments/bridge_knee_noise_eval.py
python experiments/bridge_knee_gradient_eval.py

python experiments/gradient_scaling_standalone.py

python figures/figure1_gradient_scaling.py
python figures/figure2_noise_sweep.py
python figures/figure3_fidelity_heatmap.py
\end{verbatim}

\noindent
\textbf{PDFs for this section} (relative to repo root):\
\texttt{figures/gradient\_scaling.pdf},\
\texttt{figures/noise\_sweep.pdf},\
\texttt{figures/fidelity\_heatmap.pdf}.
If a figure is missing, run the corresponding script above; if the path is
wrong, use \texttt{\textbackslash renewcommand\{\textbackslash FigPath\}\{../figures/\}}
(or your PDF folder) before \texttt{\textbackslash input} of this file.
Record Python, PennyLane, and PyTorch versions in supplementary material if needed.

% ------------------------------------------------------------------
\subsection{Scope and Limitations (Experiments)}\label{ssec:exp_limits}
% ------------------------------------------------------------------

Classical simulators only; small~$n$; single classification task for surrogate
data; linear-only templates in~(B) vs.\ mixed connectivities in chromosomes;
Experiments~(A)--(C) use \textbf{different} circuit populations---no automatic
consistency unless you re-evaluate Pareto members under~(B) or~(C).
For~(A), current search runs are small-scale and should be extended with multi-seed
convergence diagnostics (e.g., hypervolume/spread traces). For~(B), fixed
parameters isolate structural robustness but do not substitute for trained-under-
noise performance. For~(C), five-point log--log fits are exploratory without
bootstrap confidence intervals.
